% ---------------------------------------------------------------------------
% Experiment History -- closing subsection, to follow the
% "Summary of Hypotheses and Verdicts" table (structure item 9).
%
% This is where the instructions for the remaining experiments live IN THE
% PAPER. Deliberately written as design constraints derived from the history,
% never as "our supervisor asked us to" -- the rationale has to survive being
% read by someone who does not know the supervision history, and has to be
% reusable verbatim in the findings paper.
%
% Every constraint below is discharged by an experiment listed in README 3.4
% (G1-G7); the mandate-to-experiment mapping is README 2.9.
% ---------------------------------------------------------------------------

\subsection{What the History Determines}
The preceding steps are, with two exceptions, falsifications of method rather than of hypotheses about
jamming, and read together they constrain the design of everything that follows more tightly than any
positive result does.
We state those constraints here because the remaining experiments are best understood as the smallest
design that satisfies all of them simultaneously.

A policy distribution must be able to represent the statistics it is asked to produce, which a
diagonal Gaussian cannot and a finite mixture trained by policy gradient does not.
A scalar frame-level reward cannot teach an input--output correlation, because the attacker would have
to sample the correct input-dependent direction by chance before it could be reinforced.
The binding constraint there is the action parameterisation and not the reward: raw in-phase and
quadrature samples give a search space in which isotropic exploration is uninformative at any reward
shaping, so the action must instead be a small set of perturbation parameters, and the learning signal
must be a gradient wherever one is available.
Penalty terms that stand in for constraints let an attacker buy a result rather than find one, which
is why the objective of Section~\ref{sec:system:objective} carries no term for power.
Finally, a detector's headline accuracy is a statement about its training distribution and not about
its resistance, so a defender must be characterised before it is attacked, and any claim of the form
``the attacker evades the detector'' is worth only as much as the characterisation behind it.

Those constraints point at a smaller model rather than a larger one, and they are what the model of
Section~\ref{sec:system:link} was reduced to satisfy.
Reducing the link until the Neyman--Pearson test can be computed rather than trained is what lets a
negative result read ``no detector can do better than this'' instead of ``the detector we happened to
train did not do better than this''.
That is a strengthening of the defender rather than a simplification of the threat model, and it is
the only configuration in which the quantity we actually want -- the distance between what an optimal
detector achieves and what a learned one achieves, which is the room an adapting attacker has to work
in -- is measurable rather than estimated.

The reporting conventions of Section~\ref{sec:system:objective} -- comparison at matched
detectability, a stealth budget set to the defender's own false-alarm rate, and the omniscient
counter-signal drawn in every results figure -- are each a correction to one of the failures above,
and are not restated here.
